%%%%%%%%%%%%%%%%%%%%%%%%%%%%%%%%%%%%%%%%%%%%%%%%%%%%%%%%%%%%%%%%%%%%%%
% Lớp 11 | Hàm số mũ và hàm số lôgarít | Hàm số mũ. Hàm số lôgarít
% Dạng 1D6?3-1: Câu hỏi lý thuyết hàm số lũy thừa, mũ, lôgarít
% Số câu: 10  (sinh tự động bởi tach_cau.py; từ nay thầy sửa trực tiếp file này)
%%%%%%%%%%%%%%%%%%%%%%%%%%%%%%%%%%%%%%%%%%%%%%%%%%%%%%%%%%%%%%%%%%%%%%

%%%%%%%%%%%%%%%% TRẮC NGHIỆM 4 PHƯƠNG ÁN %%%%%%%%%%%%%%%%

% nguồn: Ok_TracNghiemLop11_Lan1
\begin{ex}%[1D6N3-1]%[ma:6d9bc5f5]
	Trong các hàm số sau  đây hàm số nào \textbf{không phải} là hàm số mũ?
	 \choice
		 {\True$y=x^{-4}$}
		{$y=\left(  \sqrt{3}\right) ^x$}
		{$y=4^{-x}$} 
		{$y=5^{\frac{x}{3}}$} 
	\loigiai{ 
		$y=x^{-4}$  là không phải là hàm số mũ.
	 }
\end{ex}

% nguồn: Ok_TracNghiemLop11_Lan1
\begin{ex}%[1D6N3-1]%[ma:6e37254d]
	Cho $ a>0,m,n\in \mathbb{R}$. Khẳng  định nào sau đây đúng?
	\choice 
		 {${a^m}+{a^n}={a^{m+n}}$}
		{${a^m}\cdot {a^n}={a^{m-n}}$} 
		 {\True ${({a^m})^n}={({a^n})^m}$}
		 {$\dfrac{{a^m}}{{a^n}}={a^{n-m}}$}
	\loigiai{ 
	Tính chất  lũy thừa}
\end{ex}

% nguồn: Ok_TracNghiemLop11_Lan1
\begin{ex}%[1D6N3-1]%[ma:4eb5a457]
	Tập giá trị của hàm số  $y=\mathrm{e}^{-2x+4}$ là
	\choice 
		 {\True$(0;+\infty)$}
		{$[0;+\infty)$} 
		 {$\mathbb{R}\setminus\{0\}$}
		 {$\mathbb{R}$}
	 \loigiai{
		Vì $\mathrm{e}^{-2x+4}>0\ \forall x\in\mathbb{R}$ nên tập  giá trị của hàm số là $(0;+\infty)$.
	} 
\end{ex}

% nguồn: Ok_TracNghiemLop11_Lan1
\begin{ex}%[1D6N3-1]%[ma:5d05380c]
	\immini{Đồ thị dưới đây có thể là  đồ thị của hàm số nào?
	 \choice
		{$y=3^x$} 
		 {$y=\left(\sqrt{3}\right)^x$}
		{$y=\left(-\dfrac{1}{3}\right)^x$} 
		{\True $y=\left(\dfrac{1}{3}\right)^x$}}{\begin{tikzpicture}[line cap=round,line join=round,>=latex,font=\footnotesize]
		\def\xmin{-2} 
		\def\xmax{4} 
		\def\ymin{-1} 
		\def\ymax{4} 
		\draw[-stealth]  (\xmin,0)--(0,0) node[below left]{$O$}--(\xmax,0) node[right]{$x$};
		\draw[-stealth] (0,\ymin)--(0,\ymax) node[left]{$y$}; 
		 \draw[dashed] (-1,0) node[below]{$-1$}--(-1,3)--(0,3) node[right]{$3$};
		\draw[smooth,samples=200,domain=-1.2:3]  plot (\x,{(1/3)^(\x)});
	\end{tikzpicture}
		} 
	 \loigiai{
		Đồ thị của hàm số  là $y=\left(\dfrac{1}{3}\right)^x$.
	 }
\end{ex}

% nguồn: Ok_TracNghiemLop11_Lan1
\begin{ex}%[1D6N3-1]%[ma:5d53a00d]
	Hàm số nào sau  đây là hàm số mũ?
	\choice 
		{$y=x^{2 x+1}$} 
		{\True $y=3^x$} 
		 {$y=x^3$}
		{$y=3  x$}
	 \loigiai{
		Hàm số $y=3^x$ là hàm  số mũ cơ số $a=3$.
	 }
\end{ex}

% nguồn: Ok_TracNghiemLop11_Lan1
\begin{ex}%[1D6N3-1]%[ma:e8ac35e6]
	Cho $a$ là  số thực dương và khác $1$. Mệnh đề nào sau đây là \textbf{sai}?
	 \choice
		{$\log_a\left(\dfrac{x}{y}\right)=\log_a x-\log_a y,\forall  x>0,y>0$}
		{$\log_a(xy)=\log_a  x+\log_a y,\forall x>0,y>0$}
		 {\True $\log_a x^2=\dfrac{1}{2}\log_a x,\forall x>0$}
		{$\log a=\dfrac{1}{\log_a 10}$} 
	\loigiai{Với $0<a\neq 1$, $x>0$ ta có $\log_a  x^2=2\log_a x$ nên mệnh đề $\log_a x^2=\dfrac{1}{2}\log_a x,\forall x>0$ sai.}
\end{ex}

% nguồn: Ok_TracNghiemLop11_Lan1
\begin{ex}%[1D6H3-1]%[ma:0eb4860e]
	Với $a$ là số thực dương khác $1$ tuỳ  ý, giá trị $\log_{a^2}a^3$ bằng
	\choice 
		 {$8$}
		 {$6$}
		 {$\dfrac{2}{3}$}
		{\True  $\dfrac{3}{2}$}
	\loigiai{ 
	Với $a$ là số thực dương khác $1$, ta có $\log_{a^2}a^3=\dfrac{3}{2}\log_aa=\dfrac{3}{2}$.} 
\end{ex}

% nguồn: Ok_TracNghiemLop11_Lan1
\begin{ex}%[1D6H3-1]%[ma:bcc81fc4]
	Cho hai số dương $a$, $b$, $a \neq 1$, thỏa  mãn $\log _{a^2} b+\log _a b^2=2$. Tính $\log _a b$.
	\choice 
		 {$\dfrac{8}{5}$}
		 {\True $\dfrac{4}{5}$}
		 {$2$}
		 {$4$}
	 \loigiai{
		Với $a$, $b$ là hai số dương, $a  \neq 1$, ta có \allowdisplaybreaks
		\begin{eqnarray*}
			&&  \log _{a^2} b+\log _a b^2=2 \\
			&\Leftrightarrow& \dfrac{1}{2}\log  _a b+2\log _a b=2 \\
			 &\Leftrightarrow& \dfrac{5}{2}\log _a b=2\\
			&\Leftrightarrow&  \log _a b=\dfrac{4}{5}.
		\end{eqnarray*}
	} 
\end{ex}


%%%%%%%%%%%%%%%% ĐÚNG/SAI %%%%%%%%%%%%%%%%

% nguồn: Ok_dungsai_lop11_Lan2
\begin{ex}%[1D6H3-1]%[ma:17bda1e7]
	Cho hàm số $y=2^x$. 
	 \choiceTF
		{Tập xác định  của hàm số là khoảng $(0;+\infty)$}
		 {\True Đồ thị hàm số luôn đi qua điểm $(0;1)$}
		{\True Hàm số  đồng biến trên $\mathbb{Z}$}
		{Đồ thị hàm số $y=2^x$ cắt đường thẳng $y=\dfrac{1}{8}$ tại  điểm $M(-3;8)$}
	\loigiai{ 
		\begin{itemchoice}
			\itemch Hàm số mũ $y=a^x$ có  tập xác định là $\mathbb{Z}$.
			\itemch Ta có $y(0)=2^0=1$, vậy đồ  thị luôn đi qua điểm $(0;1)$.
			\itemch Cơ số $a=2>1$  nên hàm số đồng biến trên $\mathbb{Z}$.
			\itemch Phương trình hoành  độ giao điểm
			\[2^x=\dfrac{1}{8}\Leftrightarrow 2^x=2^{-3}\Leftrightarrow  x=-3.\]
			Vậy tọa độ giao điểm của hai đồ thị  là $M\left(-3; \dfrac{1}{8}\right)$.
		\end{itemchoice}
	 }
\end{ex}


%%%%%%%%%%%%%%%% TỰ LUẬN %%%%%%%%%%%%%%%%

% nguồn: Ok_TracNghiemLop11_Lan1
\begin{ex}%[1D6N3-1]%[ma:826cca16]
	Với $a$ là số thực  dương bất kì, mệnh đề nào dưới đây \textbf{đúng}?
	 
	{$\log3a=\dfrac{1}{3}\log  a$}
	 {\True $\log a^3=3\log a$}
	{$\log(3a)=3\log a$} 
	{$\log a^3=\dfrac{1}{3}\log a$} 
	\loigiai{ 
		Mệnh đề đúng là  $\log a^3=3\log a$.
	 }
\end{ex}
